\begin{pastexam}[label={px:ncku-113-prog-07}]{成大資訊聯招 113 學年度 第 7 題}
（單選）

\textbf{[5\%]} Given an initial empty Fibonacci heap (F heap). After a sequence of insertion
operations and delete-min operations, the Fibonacci heap becomes as follows:
\begin{center}
\begin{tikzpicture}[x=0.0145cm,y=-0.0145cm,wn/.style={ellipse,draw,fill=white,inner sep=0pt,minimum width=0.86cm,minimum height=0.56cm,font=\sffamily\bfseries\small},bn/.style={wn,fill=black,text=white}]
\node[wn] (n8) at (474,281) {8};
\node[wn] (n10) at (292,352) {10};
\node[wn] (n15) at (387,352) {15};
\node[wn] (n21) at (474,352) {21};
\node[wn] (n14) at (205,424) {14};
\node[wn] (n12) at (292,424) {12};
\node[wn] (n18) at (387,424) {18};
\node[wn] (n30) at (205,507) {30};
\draw (n8)--(n10);
\draw (n8)--(n15);
\draw (n8)--(n21);
\draw (n10)--(n14);
\draw (n10)--(n12);
\draw (n15)--(n18);
\draw (n14)--(n30);
\end{tikzpicture}
\end{center}

Then, please proceed to perform operations in the following order: decrease the key 14 by 5,
decrease the key 21 by 14, delete the key 12, insert the key 14, delete the minimum key.
Considering these operations, which of the following statements is true?
\begin{enumerate}[label=(\Alph*)]
\item After performing ``delete the key 12'', the F heap has two min trees.
\item A cascading cut occurs when performing ``decrease the key 14 by 5.''
\item The resultant F heap has 4 min trees.
\item The resultant F heap has 2 min trees.
\item None of the above.
\end{enumerate}
\end{pastexam}

\begin{pastsoln}
答案是 (D)。

\textbf{約定}。 圖上沒有標記任何節點，所以一開始每個節點的標記（ChildCut）都是 false。 題目用的「F heap」
「min tree」是 Horowitz 課本的說法，那裡的「刪除任意節點」是：把節點從所在的串列拿掉，它的子節點
併入根串列，再從它的父節點開始做連鎖剪下，\textbf{不}整理根串列。 資結書第~8.6~節（和 CLRS）的
刪除則是「減少鍵值到 $-\infty$，再取出最小值」，取出最小值時會整理根串列。 兩種做法在 (A) 會得到
不同的結果，下面兩種都寫出來；最後的結果兩種一樣。 下面用 $x(\ldots)$ 表示以 $x$ 為根、括號裡是
子樹，$^{*}$ 表示有標記。

\textbf{1. 把 14 減 5，變成 9}。 $9 < $ 父節點 $10$，違反堆積性質，把 $9$ 連同子樹 $30$ 剪下放進根
串列。 父節點 $10$ 原本沒有標記，而且不是根，做標記，停止。 \textbf{沒有}連鎖剪下。
\[
8\bigl(10^{*}(12),\ 15(18),\ 21\bigr),\quad 9(30).
\]
\textbf{2. 把 21 減 14，變成 7}。 $7 < 8$，剪下放進根串列；父節點 $8$ 是根，不做標記。 最小值變成 $7$。
\[
8\bigl(10^{*}(12),\ 15(18)\bigr),\quad 9(30),\quad 7.
\]
\textbf{3. 刪除 12}。 $12$ 沒有子節點。 把它從 $10$ 的子節點串列拿掉之後，從父節點 $10$ 開始連鎖剪下：
$10$ 已經有標記，剪下放進根串列並清除標記；再往上是根 $8$，停止。
\begin{itemize}
\item Horowitz 的做法到這裡結束，有 $4$ 棵最小樹：$8(15(18))$、$9(30)$、$7$、$10$。
\item 資結書的做法是把 $12$ 減到 $-\infty$（同樣剪下 $12$、連鎖剪下 $10$），再取出最小值 $-\infty$，
接著整理根串列：$8$、$9$ 都是 $1$ 階，連成 $8(15(18), 9(30))$；$7$、$10$ 都是 $0$ 階，連成 $7(10)$。
剩 $2$ 棵。
\end{itemize}
\textbf{4. 插入 14}：直接加進根串列。

\textbf{5. 取出最小值}。 最小值是 $7$。
\begin{itemize}
\item Horowitz：根串列是 $8(15(18))$、$9(30)$、$7$、$10$、$14$。 拿掉 $7$（沒有子節點），整理：$10$、
$14$ 都是 $0$ 階，連成 $10(14)$；這時 $8$、$9$、$10$ 都是 $1$ 階，$8$ 和 $9$ 連成 $8(15(18), 9(30))$，
是 $2$ 階；剩下 $10(14)$ 是 $1$ 階。
\item 資結書：根串列是 $8(15(18), 9(30))$、$7(10)$、$14$。 拿掉 $7$，它的子節點 $10$ 進根串列；整理：
$14$、$10$ 都是 $0$ 階，連成 $10(14)$。
\end{itemize}
兩種做法最後都是兩棵最小樹：
\begin{center}
\begin{tikzpicture}[x=0.9cm,y=0.9cm,
  fn/.style={ellipse, draw, minimum width=0.86cm, minimum height=0.56cm, inner sep=0pt, font=\sffamily\small},
  fe/.style={draw=black!65}]
\node[fn] (r8) at (1,0) {8};
\node[fn] (r15) at (0,-1) {15};
\node[fn] (r9) at (2,-1) {9};
\node[fn] (r18) at (0,-2) {18};
\node[fn] (r30) at (2,-2) {30};
\node[fn] (r10) at (4,0) {10};
\node[fn] (r14) at (4,-1) {14};
\draw[fe] (r8)--(r15) (r8)--(r9) (r15)--(r18) (r9)--(r30) (r10)--(r14);
\draw[fe, densely dashed] (r8)--(r10);
\end{tikzpicture}
\end{center}
（虛線表示兩個根在同一個根串列裡；最小指標指向 $8$。）

逐項檢查：
\begin{itemize}
\item (A) 照 Horowitz 的刪除，刪掉 $12$ 之後有 $4$ 棵最小樹，錯。 （照資結書「減到 $-\infty$ 再取出
最小值」的刪除會有 $2$ 棵；但那樣 (A)、(D) 都對，和單選不合，所以題目用的是 Horowitz 的刪除。）
\item (B) 錯。 減少 $14$ 時父節點 $10$ 原本沒有標記，只是做上標記，沒有連鎖剪下。 連鎖剪下發生在
刪除 $12$ 的時候（剪下 $10$）。
\item (C) 錯，最後是 $2$ 棵。
\item (D) 對。
\end{itemize}
\end{pastsoln}
